\section{Background}
\label{sec:background}

This section introduces the concepts required to understand the experimental work.

\subsection{The RPL Routing Protocol}
\label{subsec:rpl_protocol}

The \gls{rpl} is a distance-vector routing protocol designed for resource-constrained \glspl{lln}. 
\gls{rpl} organizes nodes into a \gls{dodag}, whose topology is rooted at a node providing connectivity towards an external network.

\subsubsection{DODAG Construction and Topology Management}

Each node in a \gls{dodag} is assigned a \textit{Rank}, which represents its position relative to the root. 
The root has the lowest Rank, while nodes farther from the root generally have higher Rank values. 
Rank is used to maintain the directed acyclic structure of the \gls{dodag} and to support loop avoidance.

The selection of a preferred parent and the computation of the resulting Rank are determined by an \textit{\gls{of}}, which evaluates the routing cost associated with neighboring nodes. 
Two \glspl{of} commonly associated with \gls{rpl} are:

\begin{itemize}
    \item \textbf{\gls{of0}:} primarily considers the number of hops towards the root when selecting a parent;
    \item \textbf{\gls{mrhof}:} considers link-quality metrics such as the \gls{etx} to compute the path cost and applies hysteresis to limit unnecessary parent changes.
\end{itemize}

In the RPL implementation considered in this work, \gls{mrhof} is used for parent selection. 
The advertised Rank of a neighbor therefore represents an important component of the routing cost considered when evaluating potential parents.

\subsubsection{RPL Control Messages}

\gls{rpl} uses \gls{icmpv6} control messages to construct and maintain the \gls{dodag}:

\begin{itemize}
    \item \textbf{\gls{dio}:} advertises \gls{dodag} information, including the current Rank, version, and other configuration parameters. Nodes use received DIOs to discover potential parents;
    
    \item \textbf{\gls{dis}:} requests \gls{dio} messages from neighboring nodes, for example when joining or recovering from a disconnected state;
    
    \item \textbf{\gls{dao}:} advertises routing information towards the root, allowing downward routes to be maintained;
    
    \item \textbf{\gls{dao-ack}:} optionally acknowledges the reception of a \gls{dao}.
\end{itemize}

Among these messages, the \gls{dio} is particularly relevant to the attacks considered in this work, since it contains the Rank advertised by a node and is therefore involved in parent selection.

\subsection{Routing Attacks in RPL Networks}
\label{subsec:routing_attacks}
The distributed nature of \gls{rpl} allows a compromised internal node to participate in the routing protocol while providing malicious routing information. 
Two routing attacks are considered in this project: Blackhole and Sinkhole.

\subsubsection{Blackhole Attack}
A Blackhole attack consists of a compromised routing node that discards packets that it is expected to forward instead of delivering them towards their destination. 
The malicious node may continue to participate normally in \gls{rpl} topology formation, while disrupting data traffic that is routed through it.

The attack primarily affects packet forwarding and can therefore result in packet loss and reduced network delivery performance.

\subsubsection{Sinkhole Attack}
A Sinkhole attack attempts to attract network traffic by making a compromised node appear to provide a highly favorable route towards the \gls{dodag} root. 
In \gls{rpl}, this can be achieved by manipulating the Rank advertised in \gls{dio} messages.

The attacker advertises a Rank lower than its actual Rank, causing neighboring nodes to potentially consider it a more attractive parent. 
Once selected as a preferred parent, the malicious node can attract traffic originating from other nodes in its neighborhood.

The effectiveness of the attack depends on the routing \gls{of}, link metrics, and parent-selection mechanisms. 
In particular, the hysteresis mechanism of \gls{mrhof} can delay parent changes when the new route does not provide a sufficiently large improvement.

A Sinkhole attack can also be combined with a Blackhole attack, where the Sinkhole redirects traffic towards the compromised node and the Blackhole subsequently discards it.

\subsection{Intrusion Detection Systems (IDS) in IoT}
\label{subsec:ids_background}
\gls{ids} provide an additional layer of defense by monitoring network activity and identifying behavior that deviates from expected protocol operation.

IDS approaches can be deployed directly on network nodes or at a centralized monitoring point. 
They can also rely on different detection strategies, such as predefined protocol rules or the identification of anomalous behavior.

In this work, the IDS is based on RPL routing information and focuses on identifying inconsistencies in the observed routing hierarchy. 
In particular, Rank information advertised through \gls{dio} messages and parent relationships reported through \gls{dao} messages can be used to assess whether the resulting topology is consistent with the expected RPL hierarchy.